% Finite transfer theorems; included by the main manuscript.
% This file deliberately has no preamble, and can be \input into the main report.
% All prime-number information is isolated in hypothesis (H_theta).

\section{A finite transfer theorem for the logarithmic integral}
\label{sec:finite-transfer}

Put
\[
 \vartheta(n)=\sum_{\substack{p\le n\\p\text{ prime}}}\log p,
 \qquad \pi(n)=\#\{p\le n:p\text{ prime}\},
 \qquad \operatorname{Li}_2(x)=\int_2^x\frac{dt}{\log t}.
\]
The lower limit $2$ removes the singularity at $1$.  Changing to the
principal-value convention $\operatorname{li}(x)$ changes the answer by a
fixed computable constant; it changes none of the error classes below.
All statements about $n$ in this section concern integers $n\ge2$.
The integrals and logarithms are abbreviations for their explicit rational
approximations, as described in the foundational section.

\subsection{The exact finite identity}

For any finite sequence $a_2,\ldots,a_n$, let
$A_k=\sum_{j=2}^k a_j$.  For any $b_2,\ldots,b_n$, distributivity and
telescoping give
\begin{equation}\label{eq:finite-abel}
 \sum_{j=2}^n a_jb_j
   =A_nb_n+\sum_{k=2}^{n-1}A_k(b_k-b_{k+1}).
\end{equation}
Indeed, the coefficient of $a_j$ on the right is
$b_n+\sum_{k=j}^{n-1}(b_k-b_{k+1})=b_j$.
This is an identity of finite sums over an arbitrary commutative ring;
its proof uses only induction on the length of the sum.

Set $a_j=\log j$ when $j$ is prime and $a_j=0$ otherwise, and set
$b_j=1/\log j$.  Write
\[
 E_k=\vartheta(k)-k,\qquad
 \Delta_k=\frac1{\log k}-\frac1{\log(k+1)},\qquad
 J_n=\frac n{\log n}+\sum_{k=2}^{n-1}k\Delta_k.
\]
Then \eqref{eq:finite-abel} is exactly
\begin{equation}\label{eq:pi-discrete-transfer}
 \pi(n)-J_n=\frac{E_n}{\log n}
                  +\sum_{k=2}^{n-1}E_k\Delta_k.
\end{equation}
Telescoping the definition of $J_n$ also gives
\begin{equation}\label{eq:discrete-reference}
 J_n=\frac2{\log2}+\sum_{k=3}^n\frac1{\log k}.
\end{equation}
Since $t\mapsto1/\log t$ is decreasing on $[2,n]$, the right and left
Riemann sums on the unit partition imply
\[
 \sum_{k=3}^n\frac1{\log k}
 \le \operatorname{Li}_2(n)
 \le \sum_{k=2}^{n-1}\frac1{\log k}.
\]
Consequently,
\begin{equation}\label{eq:discrete-Li-error}
 0\le J_n-\operatorname{Li}_2(n)\le\frac2{\log2}\le4.
\end{equation}
The estimates $1/2\le\log2\le3/4$ suffice everywhere below.  They follow,
for example, by integrating $1/t$ on $[1,2]$, with the upper bound supplied
by the trapezoid inequality for the convex function $1/t$.

Two further bounds will be used:
\begin{equation}\label{eq:delta-and-trivial}
 0\le\Delta_k\le\frac1{k(\log k)^2},
 \qquad |E_k|\le3k\log k.
\end{equation}
For the first, use $\log(1+1/k)\le1/k$ and
$\log(k+1)\ge\log k$.
For the second, $0\le\vartheta(k)\le k\log k$ and
$k\le2k\log k$ imply
$|\vartheta(k)-k|\le\vartheta(k)+k\le3k\log k$.
No theorem about the density of primes is used here.

\subsection{Transfer of an arbitrary explicit modulus}

Suppose $T:\mathbb N\to\mathbb N$ is a primitive recursive function,
$T(r)\ge2$, and the following prime-number estimate has been established:
\begin{equation}\tag{$H_\vartheta$}\label{hyp:theta-modulus}
 k\ge T(r)\quad\Longrightarrow\quad
 |\vartheta(k)-k|\le2^{-r}k
 \qquad(r\in\mathbb N).
\end{equation}
This is a hypothesis of the abstract transfer lemma. In the application here
it is supplied by Theorem~\ref{thm:theta-modulus-established}; it is not an
extra assumption of the final theorem.

\paragraph{Theorem (explicit logarithmic-integral transfer).}
Define
\[
 T_\pi(m)=\max\left\{T(m+4)^2,\ 2^{16(m+5)}\right\}.
\]
Then
\begin{equation}\label{eq:pi-modulus}
 n\ge T_\pi(m)\quad\Longrightarrow\quad
 |\pi(n)-\operatorname{Li}_2(n)|
       \le 2^{-m}\frac n{\log n}.
\end{equation}

\paragraph{Proof.}
Let $L=\log n$, $\delta=2^{-(m+4)}$, and
$q=\lceil\sqrt n\rceil$.  The integer $q$ is obtained by bounded search
for the least $q\le n$ satisfying $n\le q^2$.
The hypothesis on $n$ gives $q\ge T(m+4)$.
For $2\le k<q$, \eqref{eq:delta-and-trivial} gives
$|E_k|\Delta_k\le3/\log k\le6$.
There are at most $\sqrt n$ such indices, so their total is at most
$6\sqrt n$.
For $q\le k<n$, hypothesis \eqref{hyp:theta-modulus} and
$\log k\ge L/2$ give
\[
 |E_k|\Delta_k\le\frac\delta{(\log k)^2}
                         \le\frac{4\delta}{L^2}.
\]
There are fewer than $n$ such indices.  Applying
\eqref{eq:pi-discrete-transfer} and \eqref{eq:discrete-Li-error},
\begin{equation}\label{eq:raw-transfer}
 |\pi(n)-\operatorname{Li}_2(n)|
 \le\delta\frac nL+4\delta\frac n{L^2}+6\sqrt n+4.
\end{equation}
We now check the numerical margins explicitly.  The condition
$n\ge2^{16(m+5)}$ implies $L\ge8(m+5)\ge40$.
For $L\ge16$, $\log L\le\sqrt L\le L/4$; hence
\[
 \frac{\sqrt n}{L}
   =\exp(L/2-\log L)
   \ge\exp(L/4)
   \ge\exp(2(m+5))\ge2^{m+5}.
\]
Also $4\le\sqrt n$ and $L\ge1$.  Thus the right side of
\eqref{eq:raw-transfer} is at most
\[
 \left(\frac5{16}+\frac7{32}\right)2^{-m}\frac nL
 =\frac{17}{32}\,2^{-m}\frac nL
 \le2^{-m}\frac nL.
\]
The unused margin $15/32$ is helpful when replacing the analytic notation
by rational enclosures.  This proves the theorem.

\paragraph{Why the elementary logarithmic inequality is legitimate.}
For completeness, $\log L\le\sqrt L$ for $L\ge1$ follows on writing
$v=\sqrt L$: the function $v-2\log v$ decreases on $[1,2]$ and increases
on $[2,\infty)$, and its value at $2$ is positive because $\log2<1$.
Equivalently, the same assertion follows from the corresponding finite
integral bounds; no compactness or choice is involved.

\subsection{A useful stretched-exponential specialization}

Suppose, for fixed positive integers $C,d$, that
\begin{equation}\label{eq:theta-power-modulus}
 n\ge2^{C(r+1)^d}\quad\Longrightarrow\quad
 |\vartheta(n)-n|\le2^{-r}n.
\end{equation}
Take $D=\max(2C,16)$.  The transfer theorem immediately gives
\begin{equation}\label{eq:pi-power-modulus}
 n\ge2^{D(m+5)^d}\quad\Longrightarrow\quad
 |\pi(n)-\operatorname{Li}_2(n)|\le2^{-m}\frac n{\log n}.
\end{equation}
For $d\ge3$ one may take the smaller $D=2C$, since
$2C(m+5)^d\ge16(m+5)$.  In particular this simplification applies to
$d=10$.

The usual analytic notation for the same conclusion is
\begin{equation}\label{eq:pi-stretched}
 |\pi(n)-\operatorname{Li}_2(n)|
 \le64\,\frac n{\log n}
       \exp\left(-\frac{(\log n)^{1/d}}{2D}\right)
 \qquad(n\ge2^{D6^d}).
\end{equation}
Here is a constructive derivation that does not select an integer by an
undecidable comparison of real numbers.  Let $\ell$ be the unique integer
with $2^\ell\le n<2^{\ell+1}$, and let $q$ be the largest integer with
$Dq^d\le\ell$.  Both are obtained by bounded integer search.  Then
$q\ge6$, and \eqref{eq:pi-power-modulus} applies with $m=q-5$.
Moreover, $\ell\ge6$ and $\log2\le3/4$ give
\[
 \log n<(\ell+1)\log2\le\tfrac34(\ell+1)\le\ell.
\]
Thus
\[
 q> (\ell/D)^{1/d}-1
      \ge(\log n)^{1/d}/D-1,
\]
where $D^{1/d}\le D$ was used.  Therefore
\[
 2^{-(q-5)}\le64\exp\bigl(-(\log2)(\log n)^{1/d}/D\bigr)
       \le64\exp\bigl(-(\log n)^{1/d}/(2D)\bigr),
\]
which proves \eqref{eq:pi-stretched}.
For the intended $d=10$ and $D=2C$, the exponential coefficient is
$1/(4C)$.
For the specific constant $D=2^{241}$ we can retain the tenth root
instead of replacing it by $D$. Since $D^{1/10}<2^{25}$ and
$\log2\ge1/2$, the same calculation gives the stronger explicit form
\begin{equation}\label{eq:fixed-stretched}
 |\pi(n)-\operatorname{Li}_2(n)|
 \le64\frac n{\log n}
       \exp\!\left(-2^{-26}(\log n)^{1/10}\right),
 \qquad n\ge2^{D6^{10}}.
\end{equation}
All thresholds derived below from the weaker generic coefficient remain
valid. This modest simplification changes no arithmetic construction.


\subsection{Every inverse-logarithmic order, with a primitive recursive threshold}

The stretched-exponential estimate implies much more than
$\pi(n)\sim\operatorname{Li}_2(n)$.  For $k\in\mathbb N$ set
\[
 M(k)=1+16D^2d^2k^2+24D,
 \qquad N(k)=2^{2M(k)^d}.
\]
Then
\begin{equation}\label{eq:all-log-orders}
 n\ge N(k)\quad\Longrightarrow\quad
 |\pi(n)-\operatorname{Li}_2(n)|
        \le\frac n{(\log n)^{k+1}}.
\end{equation}
To verify this, put $L=\log n$ and $u=L^{1/d}$.
The lower bound on $n$ gives $u\ge M(k)$.  In particular,
$u\ge24D$ and $u\ge16D^2d^2k^2$.
Since $\log u\le\sqrt u$ for $u\ge1$,
\[
 \log64+kd\log u
 \le6+kd\sqrt u
 \le\frac{u}{4D}+\frac{u}{4D}
 =\frac{u}{2D}.
\]
Exponentiating yields
$64\exp(-u/(2D))\le u^{-kd}=L^{-k}$.
The same lower bound on $n$ is stronger than $n\ge2^{D6^d}$, because
$M(k)\ge24D$.  Combining with \eqref{eq:pi-stretched} proves
\eqref{eq:all-log-orders}.
When $d=10$ and $D=2C$, this is the explicit formula
\[
 M(k)=1+6400C^2k^2+48C.
\]
The thresholds are intentionally coarse.  Their purpose is to be
transparent, uniform primitive recursive witnesses.

\subsection{The full factorial asymptotic expansion}

For an integer $r\ge1$ define
\[
 S_r(n)=n\sum_{j=0}^{r-1}\frac{j!}{(\log n)^{j+1}},
 \qquad B_r=1+r!\bigl(1+2^{r+2}\bigr).
\]
Repeated integration by parts, with every boundary term retained, gives
\begin{equation}\label{eq:Li-expansion-exact}
 \operatorname{Li}_2(n)
  =S_r(n)
   -2\sum_{j=0}^{r-1}\frac{j!}{(\log2)^{j+1}}
   +r!\int_2^n\frac{dt}{(\log t)^{r+1}}.
\end{equation}
For clarity, the one-step identity behind this formula is
\[
 \int_2^n\frac{dt}{(\log t)^s}
 =\frac n{(\log n)^s}-\frac2{(\log2)^s}
       +s\int_2^n\frac{dt}{(\log t)^{s+1}}.
\]
It follows by integrating the explicit derivative of
$t/(\log t)^s$; induction on $s$ yields
\eqref{eq:Li-expansion-exact}.

Suppose $L=\log n\ge16(r+2)^2$.  Splitting the remaining integral at
$\sqrt n$ gives
\[
 \int_2^n\frac{dt}{(\log t)^{r+1}}
 \le 2^{r+1}\sqrt n+\frac{2^{r+1}n}{L^{r+1}}
 \le\bigl(1+2^{r+1}\bigr)\frac n{L^{r+1}}.
\]
For the second inequality it is enough that
$(r+1)(\log2+\log L)\le L/2$.
Writing $v=\sqrt L\ge4(r+2)$, the left side is at most
$(r+1)(1+v)\le(v/4)(2v)=L/2$.
The boundary constant satisfies
\[
 2\sum_{j=0}^{r-1}\frac{j!}{(\log2)^{j+1}}
 \le2\sum_{j=0}^{r-1}j!2^{j+1}
 \le2^{r+1}r!.
\]
Also $n/L^{r+1}\ge1$ under the same lower bound on $L$.
It follows that
\begin{equation}\label{eq:Li-expansion-bound}
 |\operatorname{Li}_2(n)-S_r(n)|
 \le r!\bigl(1+2^{r+2}\bigr)\frac n{L^{r+1}}.
\end{equation}
Combining with \eqref{eq:all-log-orders}, we obtain the explicit theorem
\begin{equation}\label{eq:pi-expansion-bound}
 \begin{split}
 n&\ge2^{\max\{2M(r)^d,\,32(r+2)^2\}}\\
 &\Longrightarrow\quad
 \left|\pi(n)-n\sum_{j=0}^{r-1}\frac{j!}{(\log n)^{j+1}}\right|
 \le B_r\frac n{(\log n)^{r+1}}.
 \end{split}
\end{equation}
Thus every fixed number of terms of the usual factorial expansion is
valid, and both the coefficient and the threshold are given by primitive
recursive expressions in the requested order.

There are also effective one-sided statements.  Apply
\eqref{eq:pi-expansion-bound} with $r+1$ terms and additionally require
$\log n\ge2B_{r+1}/r!$.  Since
$S_{r+1}(n)-S_r(n)=r!n/(\log n)^{r+1}$, this gives
\[
 \frac{r!}{2}\frac n{(\log n)^{r+1}}
 \le\pi(n)-S_r(n)
 \le\frac{3r!}{2}\frac n{(\log n)^{r+1}}.
\]
For example, the extra natural-number condition
$n\ge2^{4B_{r+1}}$ is sufficient.  Consequently every fixed truncation
$S_r$ eventually underestimates $\pi$, and its first omitted term gives
the leading size of the error.  This assertion is about fixed
truncation order; it does not assert convergence of the factorial series.

\subsection{A constructive prime-in-an-interval consequence}

Under \eqref{eq:pi-power-modulus}, let $m\in\mathbb N$ and
$n\ge2^{D(m+9)^d}$.  Put $h=\lceil n/2^m\rceil$ and
$\delta=2^{-(m+4)}$.  Both $n$ and $n+h$ satisfy the estimate with
precision $m+4$, and $1\le h\le n$.
The logarithmic-integral increment satisfies
\[
 \operatorname{Li}_2(n+h)-\operatorname{Li}_2(n)
 \ge\frac h{\log(n+h)}\ge\frac h{2\log n}.
\]
The sum of the two endpoint error bounds is at most
$3\delta n/\log n$, because $n+h\le2n$ and
$\log(n+h)\ge\log n$.
As $h\ge n/2^m=16\delta n$, we conclude
\[
 \pi(n+h)-\pi(n)
 \ge\frac{h/2-3\delta n}{\log n}
 \ge\frac{5\delta n}{\log n}>0.
\]
Therefore a prime exists in $(n,n+\lceil n/2^m\rceil]$.
It is found by bounded search in that explicitly given interval.
This is a witness-producing consequence, not merely a negated
statement about empty intervals.  There is no hidden search for a
precision at which strict positivity becomes visible: since
$\log n\le n$, the displayed lower bound is at least $5\delta$.
A predetermined precision larger than $m+10$ therefore suffices to
certify positivity with rational enclosures.

\subsection{What the transfer does and does not prove about iPRA}

The prime-number hypothesis must itself be proved in iPRA; ordinary
computability of $T$ is not a substitute for such a proof.
The present section supplies the following concrete reduction.

\begin{enumerate}
 \item Equation \eqref{eq:finite-abel} is proved for coded finite rational
 lists by quantifier-free induction on list length.  It remains exact if
 the $\log j$ are replaced by positive rational approximations and each
 $b_j$ is defined as the reciprocal of that same approximation.
 \item Primality, $\pi(n)$, finite sums, integer square roots, binary
 length, fixed integer powers, and the thresholds displayed above are
 primitive recursive.  Bounded search is always given an explicit
 natural-number upper bound.
 \item The bounds for logarithms follow from finite rational sums and
 geometric-series tails.  Riemann-sum inequalities for $\operatorname{Li}_2$
 are finite inequalities with an explicit mesh error.  On $[2,n]$,
 $1/\log t\le2$ and $|\frac{d}{dt}(1/\log t)|\le2$ suffice for a uniform
 rational integration modulus.
 \item Each displayed real inequality can be replaced by the assertion
 that every precision-$s$ rational enclosure satisfies the desired
 inequality with an additional explicit $2^{-s}$ allowance.
 For a given $n,m,s$, all sums and all enclosures involved have
 primitive recursive bounds.  One first proves the corresponding
 bounded rational inequality by quantifier-free induction, and only
 then universally generalizes $n,m,s$.  One does not apply induction
 to a formula containing a quantifier over all approximation precisions.
 \item The integration-by-parts identity for the factorial expansion is
 established by the finite telescoping identity for the explicit
 derivatives, together with a Taylor-remainder bound for each rational
 mesh interval.  The derivative bounds involve only rational powers
 of a positive lower bound for $\log t$; the mesh can therefore be
 chosen by a primitive recursive formula.  Induction in the expansion
 order is again performed on the finite rational identity with its
 explicit error allowance.
\end{enumerate}

For a particularly simple all-arithmetic conclusion, let $b(n)$ be the
binary length of $n$, and let $A(n,s)\in\mathbb Q$ be an explicitly
constructed approximation to $\operatorname{Li}_2(n)$ with
$|A(n,s)-\operatorname{Li}_2(n)|\le2^{-s}$.
Since $b(n)\le4\log n$, \eqref{eq:pi-modulus} at precision $m+2$ implies
\[
 n\ge T_\pi(m+2)\quad\Longrightarrow\quad
 |\pi(n)-A(n,s)|
       \le\frac{n}{2^m b(n)}+2^{-s}.
\]
After clearing positive integer denominators this is a decidable
arithmetical inequality.  This is the recommended final object-language
statement for the Lean-to-iPRA development.

\paragraph{Lean dependency outline.}
First implement finite Abel summation over coded rational lists and
prove its coefficient-telescoping identity.  Next prove the prime
specialization, the discrete-reference formula, and the three bounds
\eqref{eq:discrete-Li-error}--\eqref{eq:delta-and-trivial}.  The proof of
\eqref{eq:raw-transfer} then consists only of splitting a finite sum at
the bounded integer square root.  The remaining modulus calculations
are elementary natural-number and rational inequalities.  The logical
separation should be reflected in Lean: the input theorem about
$\vartheta$ is an explicit hypothesis or an imported proved theorem;
it must not be hidden behind an axiom named as though it were already
established.  Proving the corresponding theorem using Mathlib's real
numbers is a useful reference development, but a separate derivation
or checked translation into the object theory is required before
claiming an iPRA theorem.
